\chapter{API 文档与健康检查：springdoc、Actuator}

后端能跑了，还要让别人（前端、运维、面试官）\textbf{看得懂、查得到状态}。这一关加两样东西：自动生成的 API 文档（Swagger 页面），和能看到“推理服务现在还好吗”的健康检查。

\partinfo{%
  \code{ops/ForecastServiceHealthIndicator.java}：10.1\\
  \code{ops/OpenApiConfig.java}、\code{forecast/StationForecastController.java}：10.2\\
  \code{src/main/resources/application.yml}：10.3%
}{\code{./mvnw test -Dtest='Part10*'}}

\begin{note}[先做完 Part 9]
\code{Part10ApiDocsTest} 会启动整个应用，所以 Part 7--9 必须已经全部完成（\code{FinalApplicationTest} 是绿的）。
\end{note}

\section{知识点}

\subsection{Actuator：应用自带的“体检报告”}

加上 \code{spring-boot-starter-actuator} 依赖（已经加好了）以后，应用会多出一组管理接口，最常用的是 \code{/actuator/health}：

\begin{jsoncode}
{
  "status": "UP",
  "components": {
    "db":              { "status": "UP", "details": { "database": "MySQL" } },
    "diskSpace":       { "status": "UP" },
    "forecastService": { "status": "UP", "details": { "baseUrl": "http://ml-service:8001",
                                                      "circuitBreaker": "CLOSED" } }
  }
}
\end{jsoncode}

\begin{itemize}
  \item 只要有一个组件 \code{DOWN}，整体就是 \code{DOWN}，HTTP 状态码是 503。Docker、Kubernetes、负载均衡器都靠这个接口判断“要不要把流量发给这个实例”。
  \item \code{db}、\code{diskSpace} 是 Spring Boot 自动提供的；\code{forecastService} 是你这一关要写的。
\end{itemize}

为了安全，默认只对外暴露 \code{health}，而且不显示 \code{details}。要在 \code{application.yml} 里显式打开（10.3）。

\subsection{自定义健康检查：HealthIndicator}

实现 \code{HealthIndicator} 接口、注册成 Bean，它就会出现在 \code{/actuator/health} 里。组件名是 Bean 名去掉 \code{HealthIndicator} 后缀：\code{forecastServiceHealthIndicator} → \code{forecastService}。

\begin{java}
@Component
public class PaymentGatewayHealthIndicator implements HealthIndicator {

    @Override
    public Health health() {
        if (gateway.ping()) {
            return Health.up().withDetail("region", "eu-west-1").build();
        }
        return Health.down().withDetail("error", "gateway timeout").build();
    }
}
\end{java}

\begin{note}[包名]
Spring Boot 4 里 \code{HealthIndicator}、\code{Health} 在 \code{org.springframework.boot.health.contributor}（Boot 3 是 \code{...actuate.health}，网上的老教程会写错）。
\end{note}

\textbf{健康检查本身不能抛异常}：对方连不上时要返回 \code{DOWN}，而不是让 \code{/actuator/health} 报 500。

\subsection{springdoc：从代码生成 API 文档}

加上 \code{springdoc-openapi-starter-webmvc-ui} 依赖（已经加好了）以后，springdoc 会扫描所有 Controller，自动生成 OpenAPI 格式的文档：

\begin{itemize}
  \item \code{/v3/api-docs}：机器可读的 JSON，前端可以用它生成调用代码；
  \item \code{/swagger-ui/index.html}：给人看的网页，还能直接在页面上发请求试接口。
\end{itemize}

不写任何注解也能生成，但说明会很简陋。几个常用注解（包 \code{io.swagger.v3.oas.annotations}）：

\begin{java}
@Tag(name = "Books", description = "Book catalogue")          // 类上：分组
public class BookController {

    @Operation(summary = "Find a book", description = "...")   // 方法上：接口说明
    @GetMapping("/{isbn}")
    public BookResponse get(
            @Parameter(description = "ISBN-13") @PathVariable String isbn) { ... }   // 参数说明
}
\end{java}

文档整体的标题和版本，用一个 \code{@Bean} 返回 \code{OpenAPI} 对象来设置（又是“第三方库的对象用 \code{@Bean} 注册”）：

\begin{java}
@Bean
public OpenAPI libraryOpenApi() {
    return new OpenAPI().info(new Info().title("Library API").version("v2"));
}
\end{java}

\subsection{测试是怎么验证的}

\begin{itemize}
  \item \code{Part10HealthIndicatorTest}：\code{@RestClientTest} + \code{MockRestServiceServer} 冒充 Python 服务的 \code{GET /health}，检查你返回的 \code{Health}。
  \item \code{Part10ApiDocsTest}：启动整个应用，请求 \code{/v3/api-docs}、\code{/swagger-ui/index.html}、\code{/actuator/health}。推理服务地址被指向一个不存在的端口，所以 \code{forecastService} 一定是 \code{DOWN}，整体返回 503。
\end{itemize}

\section{题目}

\begin{task}{10.1\quad ForecastServiceHealthIndicator}
交给 Spring 管理。构造函数里用注入的 \code{RestClient.Builder} 和 \code{properties.baseUrl()} 创建 RestClient。

\code{health()}：调用 Python 服务的 \code{GET /health}，响应转成已经写好的 \code{ServiceStatus} record：
\begin{itemize}
  \item \code{status} 是 \code{"UP"} → \code{Health.up()}，其他情况 → \code{Health.down()}；
  \item 调用出错（\code{RestClientException}）→ \code{Health.down()}，加一个 detail \code{"error"}（异常消息）；
  \item 不管哪种情况，都加上 \code{"baseUrl"} 和 \code{"circuitBreaker"}（熔断器当前状态的名字，比如 \code{"CLOSED"}）两个 detail。
\end{itemize}

\begin{example}
\begin{plaincode}
GET /health -> {"status":"UP"}    ->  UP,   {baseUrl=http://..., circuitBreaker=CLOSED}
GET /health -> {"status":"DOWN"}  ->  DOWN
GET /health -> 503 / 连不上        ->  DOWN, {error=..., baseUrl=..., circuitBreaker=CLOSED}
\end{plaincode}
\end{example}
\end{task}

\begin{task}{10.2\quad API 文档}
\textbf{a)} \code{OpenApiConfig} 是配置类，\code{dublinBikesOpenApi()} 注册成 Bean，返回的 \code{OpenAPI} 标题 \code{"Dublin Bikes API"}、版本 \code{"v1"}，描述随意。

\textbf{b)} 给 Part 9 写的 \code{StationForecastController} 加文档注解：类上 \code{@Tag(name = "Forecast")}；预测方法上 \code{@Operation(summary = "Forecast bikes for the next hours")}；\code{hours} 参数上 \code{@Parameter} 写一句说明。

做完以后启动应用，打开 \url{http://localhost:8080/swagger-ui/index.html} 看看。
\end{task}

\begin{task}{10.3\quad 打开 Actuator 的详细信息}
在 \code{application.yml} 里（和 \code{spring:}、\code{app:} 同一级）加一段 \code{management} 配置：
\begin{itemize}
  \item 对外暴露 \code{health} 和 \code{info} 两个接口（\code{management.endpoints.web.exposure.include}）；
  \item \code{health} 总是显示 details（\code{management.endpoint.health.show-details: always}）。
\end{itemize}
\end{task}

\section{我直接做好的：Docker 和 CI}

这两样不属于 Spring 的写法，我已经做好放进仓库了，你读懂、跑一遍就行。

\textbf{一条命令启动整个系统：}

\begin{shell}
docker compose up --build        # MySQL + Python 推理服务 + Spring Boot，打开 http://localhost:8080
\end{shell}

\begin{itemize}
  \item \code{Dockerfile}：多阶段构建。第一阶段用 JDK 编译打包，第二阶段只放 JRE 和 jar，镜像小很多；先单独下载依赖，改代码时这一层会被缓存。
  \item \code{ml-service/Dockerfile}：模型文件约 2GB，不打进镜像，运行时从 \code{./ml-service/models} 挂载进去。
  \item \code{compose.yaml}：三个服务都有健康检查，Spring Boot 等 MySQL 和推理服务都健康了才启动；它用 \code{mysql} profile 连 MySQL，这也是 Flyway 第一次在真实的 MySQL 上建表（我测过，没有问题）。
  \item \code{.github/workflows/ci.yml}：每次 push 到 main，GitHub Actions 自动跑全部测试，再分别构建两个镜像。
\end{itemize}

\begin{think}
\begin{enumerate}
  \item \code{/actuator/health} 返回 503 时，负载均衡器会怎么做？推理服务挂了，整个后端都被判定成“不健康”合理吗？（提示：liveness 和 readiness 的区别）
  \item 为什么 Actuator 默认不暴露所有接口、不显示 details？
  \item API 文档从代码生成，和手写一份文档相比，好处和局限是什么？
  \item Docker 多阶段构建为什么能让镜像变小？为什么要先 \code{COPY pom.xml} 下载依赖，再 \code{COPY src}？
\end{enumerate}
\end{think}
